\documentclass[12pt,journal,compsoc]{article}
%PREAMBLE
\usepackage[vmargin=4cm,tmargin=3cm,hmargin=2.5cm,letterpaper]{geometry}
\usepackage[english,spanish]{babel}
\usepackage[utf8]{inputenc}
\usepackage{fancyhdr}
\usepackage{graphicx}
\usepackage{amsmath}
\usepackage{hyperref}
\usepackage{mdframed}								
\def\changemargin#1#2{\list{}{\rightmargin#2\leftmargin#1}\item[]}
\let\endchangemargin=\endlist
\setlength\parindent{0pt}
\graphicspath{{../multimedia/images/}}
%BEGIN DOCUMENT								
\begin{document}
\fancyhead[L]{\includegraphics[scale=0.15]{ucr}}
\fancyhead[C]{UNIVERSIDAD DE COSTA RICA\\ESCUELA DE INGENIERÍA ELECTRICA\\PROGRAMACIÓN BAJO PLATAFORMAS ABIERTAS\\PROYECTO DE DESARROLLO - PROPUESTA\\EQUIPO CERBERUS}
\fancyhead[R]{\includegraphics[scale=0.4]{eie}}
\pagestyle{fancy}
\setlength{\headheight}{37pt}
\rule{\linewidth}{0.5mm}
\begin{center}
\begin{LARGE}
\textbf{\textit{Implementación de MySQL para crear una Base de Datos de Notas para la Escuela}}\\
\end{LARGE}
\vspace{0.25cm}
\end{center}
MySQL es un sistema de gestión de bases de datos creado por ORACLE que utiliza SQL, lenguaje de consulta estructurado, que se define como un lenguaje declarativo de acceso a bases de datos relacionales que permite especificar diversos tipos de operaciones en ellas. Una de sus características es el manejo del álgebra y el cálculo relacional que permiten efectuar consultas con el fin de recuperar de forma sencilla información de interés de bases de datos, así como hacer cambios en ella.\\\\
Dado los beneficios que está plataforma nos brinda, la idea es implementar una base de datos para la escuela, donde está cuente con datos de los estudiantes y sus notas, y así poder obtener información como por ejemplo promedios y brindar un reporte de lo consultado.
\subsection*{Objetivo General}
\begin{itemize}
\item Adquirir las habilidades y conocimientos de MySQL necesarios para utilizar y administrar la base de datos de notas de estudiantes egresados para la EIE.
\end{itemize}
\subsection*{Objetivos Específicos}
\begin{enumerate}
\item 
\item 
\item 
\end{enumerate}
\begin{thebibliography}{6}
\bibitem{label} Silberschatz, A., Korth, H., Sudarshan, S. (2002). Fundamentos de Bases de Datos (4ta edición), McGraw Hill.
\bibitem{label}ref
\bibitem{label}ref
\bibitem{}\url{http://dev.mysql.com/doc/refman/5.0/es/tutorial.htm}
\bibitem{ros}\url{http://}
\bibitem{}\url{http://}
\end{thebibliography}
\end{document}